\documentclass[11pt]{article}
\usepackage[margin=1in]{geometry}
\usepackage{amsmath,amsthm,amssymb}
\usepackage{booktabs,hyperref,enumitem,array,xcolor,colortbl,mdframed}

\definecolor{proved}{HTML}{1a7a1a}
\definecolor{open}{HTML}{b30000}
\definecolor{withdrawn}{HTML}{555555}
\definecolor{lightgreen}{HTML}{e6f4e6}
\definecolor{lightred}{HTML}{fde8e8}
\definecolor{lightgray}{HTML}{f5f5f5}
\definecolor{darkblue}{HTML}{003366}

\newtheorem{theorem}{Theorem}[section]
\newtheorem{lemma}[theorem]{Lemma}
\newtheorem{proposition}[theorem]{Proposition}
\newtheorem{definition}[theorem]{Definition}
\newtheorem{remark}[theorem]{Remark}
\newtheorem*{openproblem}{Open Problem}

\newcommand{\Qt}{\widetilde{Q}}
\newcommand{\lmin}{\lambda_{\min}}

\title{\textbf{The Inverse-GCD Operator $Q_N$ (Q6):\\
Definitions, a Restricted Rayleigh Bound,\\
and an Exploratory Spectral Programme}\\[6pt]
\large\textit{Crown-jewel operator note --- what is proved and what is open}}
\author{Jonathan Robert Simons\\
Savannah, Georgia, USA\\
\texttt{simonsmedicalinnovations@gmail.com}}
\date{Corrected preprint \quad|\quad August 2026}

\begin{document}
\maketitle

\begin{center}
\fbox{\parbox{0.92\linewidth}{\centering\small
\textbf{Scope.}
Q6 is developed here as a spectral operator on $\{1,\ldots,N\}$.
Proved: Bridge$^*$ on single prime pairs, and the nonnegative cone.
Exploratory / open: multi-representation bounds.
False as stated: universal full-spectrum floors $\lmin>-1/2$.
No Navier--Stokes claim.}}
\end{center}
\vspace{0.6em}

\begin{mdframed}[linecolor=darkblue,linewidth=1.2pt]
{\large\textbf{\color{darkblue}Status}}\\[4pt]
\begin{tabular}{>{\raggedright\arraybackslash}p{9.2cm}
                >{\centering\arraybackslash}p{3.2cm}}
\toprule
\textbf{Claim} & \textbf{Status}\\
\midrule
\rowcolor{lightgreen}
Definitions of $Q_N$, $\Qt_N$, $H_N$ & \textcolor{proved}{\textbf{Locked}}\\
\rowcolor{lightgreen}
Bridge$^*$: $R(e_p-e_q)>-1/2$ on $\Qt_N$ & \textcolor{proved}{\textbf{Proved}}\\
\rowcolor{lightgreen}
Positive orthant: $v\ge 0\Rightarrow v^\top\Qt_N v\ge 0$ & \textcolor{proved}{\textbf{Proved}}\\
\rowcolor{lightred}
$\lmin(Q_N)>-1/2$ for all $N$ & \textcolor{withdrawn}{\textbf{False}}\\
\rowcolor{lightred}
$\lmin(\Qt_N)>-1/2$ for all $N$ & \textcolor{withdrawn}{\textbf{False}}\\
\rowcolor{lightred}
Dark-state $\Leftrightarrow$ Goldbach (June 2026 draft) & \textcolor{withdrawn}{\textbf{Withdrawn}}\\
\rowcolor{lightgray}
Multi-representation Bridge$^*$ & \textcolor{open}{\textbf{Open}}\\
\bottomrule
\end{tabular}
\end{mdframed}

\begin{abstract}
We give a corrected account of the inverse-GCD operators used in earlier
``Q6'' drafts. Write
\[
Q_N(i,j)=\frac{1}{\gcd(i,j)},\qquad
\Qt_N(i,j)=\frac{1}{\gcd(i,j)\sqrt{ij}},
\qquad 1\le i,j\le N,
\]
and let $H_N$ be the degree-normalized form of $\Qt_N$.
We prove that for distinct primes $p,q$,
\[
R(e_p-e_q)
=\frac12\Big(\frac1{p^2}+\frac1{q^2}\Big)-\frac1{\sqrt{pq}}
\;>\;-\frac12
\]
on $\Qt_N$ (Bridge$^*$).
We also record that every nonnegative vector has nonnegative Rayleigh
quotient on $\Qt_N$.
We withdraw the full-spectrum Bridge claim $\lmin>-1/2$: already for
moderate $N$, $\lmin(Q_N)<-1/2$ and $\lmin(\Qt_N)<-1/2$.
Earlier dark-state / Goldbach packaging and any Navier--Stokes claim from
these matrices are withdrawn.
\end{abstract}

\section{Operators}

\begin{definition}[Raw inverse-GCD]
\[
Q_N(i,j)=\frac{1}{\gcd(i,j)},\qquad 1\le i,j\le N.
\]
\end{definition}

\begin{definition}[Normalized inverse-GCD]
\[
\Qt_N(i,j)=\frac{1}{\gcd(i,j)\sqrt{ij}}.
\]
In particular, for distinct primes $p,q\le N$,
\[
\Qt_N(p,p)=\frac1{p^2},\qquad
\Qt_N(p,q)=\frac1{\sqrt{pq}}.
\]
\end{definition}

\begin{definition}[Degree-normalized form]
Let $d_i=\sum_{j=1}^N \Qt_N(i,j)$ and $D=\mathrm{diag}(d_1,\ldots,d_N)$.
Set
\[
H_N = D^{-1/2}\,\Qt_N\,D^{-1/2}.
\]
\end{definition}

\begin{remark}[Hygiene]
These three matrices are different.
A spectral statement about one does not automatically transfer to another,
nor to any fluid commutator called ``Theorem~H'' in companion notes.
\end{remark}

\begin{remark}[Withdrawn identity]
Some June 2026 drafts asserted
$1/n=\sum_{d\mid n}\mu(d)\varphi(d)/d^2$.
This fails already at $n=2$ ($1/2$ versus $3/4$).
No argument in this note uses that identity.
\end{remark}

\section{What fails: full-spectrum floors}

\begin{proposition}[Full-spectrum Bridge is false]\label{prop:floorfails}
There exist $N$ such that
\[
\lmin(Q_N)<-1/2
\qquad\text{and}\qquad
\lmin(\Qt_N)<-1/2.
\]
In particular the universal claims
$\lmin(Q_N)>-1/2$ and $\lmin(\Qt_N)>-1/2$ for all $N$ are false.
\end{proposition}

\begin{proof}[Numeric certificate]
Direct eigendecomposition yields, for example,
$\lmin(Q_{10})\approx -1.90<-1/2$ and
$\lmin(\Qt_{20})\approx -0.505<-1/2$.
Reproducible by the accompanying script
\texttt{scripts/bridge\_floor\_verify.py}.
\end{proof}

\begin{remark}
A claimed universal lower bound $\lmin(H_N)\ge -3/14$ is likewise not
uniform in small $N$; see the companion verification suite.
We do not package any $H_N$ floor as a theorem here.
\end{remark}

\section{Bridge$^*$: single prime pair}

\begin{theorem}[Bridge$^*$, single pair]\label{thm:bridge}
Let $p\neq q$ be primes and $N\ge\max(p,q)$. For $v=e_p-e_q$,
\[
R(v)
:=\frac{v^\top\Qt_N v}{\|v\|_2^2}
=\frac12\Big(\frac1{p^2}+\frac1{q^2}\Big)-\frac1{\sqrt{pq}}
\;>\;-\frac12.
\]
\end{theorem}

\begin{proof}
Since $\|v\|_2^2=2$,
\[
v^\top\Qt_N v
=\Qt_N(p,p)+\Qt_N(q,q)-2\Qt_N(p,q)
=\frac1{p^2}+\frac1{q^2}-\frac{2}{\sqrt{pq}},
\]
hence
\[
R=\frac12\Big(\frac1{p^2}+\frac1{q^2}\Big)-\frac1{\sqrt{pq}}.
\]
Then
\[
R+\frac12
=\frac12+\frac1{2p^2}+\frac1{2q^2}-\frac1{\sqrt{pq}}
>\frac12-\frac1{\sqrt{pq}}.
\]
Distinct primes force $pq\ge 6$, so $1/\sqrt{pq}\le 1/\sqrt6<1/2$.
Therefore $R+1/2>0$.
\end{proof}

\begin{remark}
The quotient may be negative: $(p,q)=(2,3)$ gives $R\approx -0.228>-1/2$.
The bound is a \emph{restricted} floor on difference vectors of single
prime pairs, not a full-spectrum eigenvalue floor.
\end{remark}

\section{Positive orthant}

\begin{proposition}[Nonnegative cone]\label{prop:cone}
If $v\in\mathbb{R}^N$ satisfies $v_i\ge 0$ for all $i$, then
$v^\top\Qt_N v\ge 0$.
\end{proposition}

\begin{proof}
Every entry of $\Qt_N$ is positive, so the quadratic form on the
nonnegative orthant is a sum of nonnegative terms.
\end{proof}

\section{Earlier drafts --- corrected status}

Earlier public drafts pushed Q6 toward Strong Goldbach via a
dark-state lemma and a full-spectrum Bridge floor, and toward
RH / Navier--Stokes packaging. Those extensions are not supported
by the arguments as written; they are marked exploratory-failed below
and are not part of the proved core.

\begin{enumerate}[leftmargin=1.2em]
\item \textbf{Broken detector.}
The indicator difference
$v_k(j)=1_{\mathbb{P}}(j)-1_{\mathbb{P}}(k-j)$
vanishes on genuine Goldbach pairs $p+(k-p)=k$, so it does not detect
them.
\item \textbf{False sign lemma.}
On raw $Q_N$, for $v=e_p-e_q$,
$\langle v,Q_N v\rangle=1/p+1/q-2$ is typically negative
(e.g.\ $(3,7)\mapsto -1.52$).
A claimed equivalence ``dark state $\Leftrightarrow$ no Goldbach'' does
not hold as written.
\item \textbf{No NS map.}
A spectral floor on $Q_N$, $\Qt_N$, or $H_N$ is not, by itself, a
bound on $\|(u\cdot\nabla)u\|$ for Leray--Hopf solutions on
$\mathbb{T}^3$.
\end{enumerate}

Older DOIs may remain as dated history.
The public credit for this operator is
Theorem~\ref{thm:bridge} and Propositions~\ref{prop:floorfails}--\ref{prop:cone},
plus the exploratory open problem in the last section.

\section{Open problem}

\begin{openproblem}[Multi-representation Bridge$^*$]
For even $k\le N$ let
$\mathcal{R}(k)=\{p\text{ prime}:2\le p\le k/2,\ k-p\text{ prime}\}$
and
$v_k=\sum_{p\in\mathcal{R}(k)}(e_p-e_{k-p})$
(indices $\le N$).
Prove or disprove the existence of an absolute $c_0<1/2$ such that
\[
R(v_k)=\frac{v_k^\top\Qt_N v_k}{\|v_k\|_2^2}\ge -c_0
\]
for all such $k$, whenever $v_k\neq 0$.
\end{openproblem}

\begin{remark}
Numeric checks through $N=200$ remain above $-1/2$ on sampled multi-rep
vectors; that is evidence, not a proof.
\end{remark}

\section*{Reproducibility}

\begin{verbatim}
python3 scripts/bridge_floor_verify.py 50
python3 -m unittest tests.test_bridge_star_h_n -v
\end{verbatim}

\section*{Acknowledgments}
This note supersedes the earlier Q6 preprint
(\href{https://doi.org/10.5281/zenodo.20405589}{10.5281/zenodo.20405589})
as the author's intended public write-up of the inverse-GCD operator.

\begin{thebibliography}{9}
\bibitem{ZenodoQ6}
J.R.~Simons,
\textit{A Q6 Spectral Route to the Strong Goldbach Conjecture}
(withdrawn as a closure claim; archive only),
Zenodo \href{https://doi.org/10.5281/zenodo.20405589}{10.5281/zenodo.20405589}.

\bibitem{BridgeStarMD}
J.R.~Simons,
\textit{Bridge$^*$ proof draft},
repository note \texttt{docs/BRIDGE-STAR-PROOF.md} (2026).
\end{thebibliography}

\end{document}
